\documentclass[11pt]{article}
\usepackage[margin=1in]{geometry}
\usepackage{amsmath,amssymb,amsthm,xurl,hyperref}
\hypersetup{colorlinks=true,linkcolor=blue,urlcolor=blue,
pdftitle={Sharp lower-end simplex stability from integrated witnesses and projection caps},
pdfauthor={Yongxian Zhang}}
\newtheorem{theorem}{Theorem}
\newtheorem{lemma}[theorem]{Lemma}
\newtheorem{proposition}[theorem]{Proposition}
\newcommand{\conv}{\operatorname{conv}}
\newcommand{\supp}{\operatorname{supp}}
\newcommand{\EE}{\mathbb E}
\newcommand{\pos}[1]{(#1)_+}
\newcommand{\BB}{B_2^d}
\newcommand{\vol}{\operatorname{vol}}
\title{Sharp lower-end simplex stability from\\integrated witnesses and projection caps}
\author{Yongxian Zhang\\[2pt]\small School of Computer Science and Engineering\\\small South China University of Technology\\\small\texttt{mxymmxym1@gmail.com}\\\small ORCID: 0009-0000-3864-3536}
\date{8 October 2026}
\pdfinfoomitdate=1
\pdftrailerid{}
\usepackage{xurl}
\hypersetup{pdfauthor={Yongxian Zhang}}
\begin{document}
\maketitle

\noindent\textbf{Result.}
For every fixed $d\ge3$ and every maximum-volume inscribed simplex $S$,
\[
 E(K,S)\le G_d^\sharp
 \left(a(K)-\frac1{d+1}\right)^{1/(d-1)}.
\]
All constants below are explicit. The exponent matches the exact truncation
obstruction. The additional geometric step uses centering to compare normalized
projection volumes, controls scale by first variation, and then puts a cap
in a $(d-1)$-dimensional projection. No inverse-Minkowski stability theorem
is an input. The supplied explicit $1/d$ theorem is preserved as a separate
source, with its constants unchanged. An independent analytic model audit
passed, with independently implemented finite checks. The main bound and its exponent obstruction now have complete Lean proofs;
see the formal-verification section. No external human peer review is claimed.
No priority or best-known-result claim is made.

\section{Definitions and explicit constants}

For a convex body $K\subset\mathbb R^d$, let $\Pi K$ denote the projection
body with $h_{\Pi K}(u)=|K\mid u^\perp|_{d-1}$ for unit $u$, and put
\[
 R(K)=\frac{|\Pi K|}{|K|^{d-1}},\qquad
 a(K)=\left(\frac d{d+1}\right)^d
       \frac{R(\mathcal P K)}{R(K)}-1,\qquad
 e(K)=a(K)-\frac1{d+1}.
\]
Here $\mathcal P K=\operatorname{conv}(K\times\{0\},(0,1))$ is the
pyramid over $K$; the ratio is independent of the chosen nonzero height.
For a simplex $S\subset K$ with centroid $z_S$, define
\[
 E(K,S)=\inf\{t\ge0:K\subset z_S+(1+t)(S-z_S)\}.
\]
For the constants set
\begin{align}
 R_0&=d(d+1),&
 b&=\frac1{4(dR_0)^d},&
 M&=b^{-1},\label{eq:constants1}\\
 Q&=(d+1)(d+2)8^d b^{-4d},&
 m&=d-1,&
 L&=m(M+1),\label{eq:constants2}\\
 J&=\frac d2(2R_0)^d
       \bigl[1+M+d2^{d-1}M\bigr],\label{eq:J}\\
 r_\sharp&=\min\left\{b,\,
             \frac1{J(8MdL)^m}\right\},&
 e_\sharp&=\frac{r_\sharp}{Q(d+1)},\label{eq:gate}\\
 A_d^\sharp&=4MdL[JQ(d+1)]^{1/m},&
 G_d^\sharp&=\max\{A_d^\sharp,\,
                    (R_0-1)e_\sharp^{-1/m}\}.\label{eq:global}
\end{align}
All constants other than the displayed $m$th roots are positive rationals.
They are extremely poor numerically and depend on dimension.

\begin{theorem}\label{thm:endpoint}
For every convex body $K\subset\mathbb R^d$, $d\ge3$, and every
maximum-volume inscribed simplex $S$,
\[
 E(K,S)\le G_d^\sharp e(K)^{1/(d-1)}.
\]
If $0\le e(K)\le e_\sharp$, the constant $A_d^\sharp$ applies, and
$A_d^\sharp e_\sharp^{1/(d-1)}\le1/2$.
No larger power can hold uniformly near $e=0$, even if one permits
a dimension-dependent constant.
\end{theorem}

\section{Normalization and the cone law}

Fix an arbitrary maximum simplex and map it affinely to a regular simplex
$\Delta$ with centroid zero and inradius one. Its vertices have norm $d$.
Replacing a vertex by $x\in K$ multiplies the simplex volume by the absolute
value of the corresponding barycentric coordinate. Maximality therefore
gives $|\alpha_i(x)|\le1$, and
\begin{equation}\label{eq:normalization}
 \BB\subset\Delta\subset K\subset R_0\BB.
\end{equation}
Both $a$ and $E$ are preserved by this affine normalization.

Write $h_K$ for the support function and $S_K$ for surface area measure.
Push the probability measure $h_K(u)dS_K(u)/(d|K|)$ forward by
$u\mapsto u/h_K(u)$; call the resulting law $\nu$.
The classical identities $\int u\,dS_K(u)=0$ and
$\int h_K\,dS_K=d|K|$ show that $\nu$ is centered.
It is supported on $\partial K^\circ\subset\BB$.
For independent samples from $\nu$, define
\[
 A=\EE|\det(X_1,\ldots,X_d)|,\qquad
 B=\EE|\det((X_1,1),\ldots,(X_{d+1},1))|.
\]
The entry005 cone-law representation is
\begin{equation}\label{eq:cone}
 a(K)=\frac{B}{(d+1)A},\qquad
 D:=B-A=(d+1)Ae(K)\le(d+1)e(K).
\end{equation}
Here $0<A\le1$, the upper bound following from Hadamard's inequality.
For completeness, the representation for polytopes follows by expanding
the two determinant expectations over distinct facets, with the $d!$
and $(d+1)!$ permutation factors. Polytope approximation, uniform
support-function convergence and weak convergence of surface area measures
give the arbitrary-body identity; all laws stay in a common compact set.
This representation, already established in entry005 v3, is the only
entry005-specific identity imported here.

For a body $H$ containing zero in its interior, its cone law satisfies
the Cauchy identity
\begin{equation}\label{eq:brightness}
 \int |u\cdot x|\,d\nu_H(x)
   =\frac{2|H\mid u^\perp|_{d-1}}{d|H|}
       \qquad (|u|=1).
\end{equation}
Indeed the factors $h_H$ cancel in the pushforward integral, leaving
$\int|u\cdot n|\,dS_H(n)/(d|H|)$.
Centering, Cauchy's formula, and \eqref{eq:normalization} give
\begin{equation}\label{eq:dispersion}
 \EE\pos{-u\cdot X}
 =\frac{|K\mid u^\perp|_{d-1}}{d|K|}
 \ge\frac{(2/d)^{d-1}}{d(2R_0)^d}=2b.
\end{equation}
The projected unit ball contains a cube of side $2/d$.

\section{The integrated assignment bound, including nonatomic laws}

This section reproduces the probabilistic argument in the supplied readable
fallback so that the quantitative proof is complete.
For an ordered base tuple $\mathbf x=(x_1,\ldots,x_d)$ put
\[
 F_{\mathbf x}(y)
   =\det((x_1,1),\ldots,(x_d,1),(y,1)),
\quad
 P(\mathbf x)=\EE\pos{F_{\mathbf x}(X)},
\quad
 N(\mathbf x)=\EE\pos{-F_{\mathbf x}(X)}.
\]
Since $\EE X=0$, $\EE F_{\mathbf x}(X)$ is the horizontal determinant,
up to sign. Subtracting its absolute value from $\EE|F_{\mathbf x}(X)|$
and integrating gives
\begin{equation}\label{eq:cancellation}
 D=2\EE_{\mathbf X}\min\{P(\mathbf X),N(\mathbf X)\}\ge0.
\end{equation}
No nonsingularity assumption on a sampled base is needed.

Use the continuous nonnegative witness
\[
 \psi_{\mathbf x}(y,z)
 =\min\{\pos{F_{\mathbf x}(y)},\pos{-F_{\mathbf x}(z)}\}
  +\min\{\pos{-F_{\mathbf x}(y)},\pos{F_{\mathbf x}(z)}\}.
\]
For independent $Y,Z$ each of its minima has expectation at most
$\min(P,N)$; hence
\begin{equation}\label{eq:witness}
 \EE_{\mathbf X,Y,Z}\psi_{\mathbf X}(Y,Z)\le D.
\end{equation}

\begin{lemma}[Integrated assignment]\label{lem:assignment}
For any centered law on $\BB$ satisfying \eqref{eq:dispersion}, there are
affinely independent $w_0,\ldots,w_d\in\supp\nu$ and a measurable assignment
$X\mapsto w_{I(X)}$ with
\[
 h:=\EE\|X-w_{I(X)}\|\le QD.
\]
If $QD\le b$, then $b\BB\subset T:=\conv(w_0,\ldots,w_d)\subset\BB$.
\end{lemma}

\begin{proof}
Let $W_0,\ldots,W_d$ be independent samples and let
$V(W)=|\det((W_0,1),\ldots,(W_d,1))|$.
For $\Sigma=\EE XX^{\mathsf T}$, centering and
\eqref{eq:dispersion} give $\EE|u\cdot X|\ge4b$ and
$u^{\mathsf T}\Sigma u\ge16b^2$.
Expanding determinants and using independence yields
\[
 \EE V^2=(d+1)!\det\Sigma
       \ge(d+1)!16^d b^{2d}\ge2b^{2d}.
\]
The lifted second-moment matrix is $\operatorname{diag}(\Sigma,1)$;
the determinant expansion gives $(d+1)!$ times its determinant.
Subtracting one lifted column from the others and applying Hadamard gives
$V\le2^d$. Consequently, with
\[
 v_0=b^d,\qquad q_0=b^{2d}/4^d,\qquad
 \mathcal E=\{W:V(W)\ge v_0\},
\]
we have $\mathbb P(\mathcal E)\ge q_0$, since
$\EE V^2\le b^{2d}+4^d\mathbb P(\mathcal E)$.

For each $i$ use the other $d$ anchors as base and test $w_i,x$.
For each pair $i<j$, use $x$ and the other $d-1$ anchors as base and
test $w_i,w_j$.
Let $H(W)$ be the integral in $x$ of their sum; there are
$N_d=(d+1)(d+2)/2$ summands.
Each uses $d+2$ independent samples before conditioning, so
\eqref{eq:witness} gives $\EE H(W)\le N_dD$.
Conditioning on $\mathcal E$, choose anchors with
$H(W)\le N_dD/q_0$.
This selection is valid for nonatomic laws: $H$ is continuous by dominated
convergence, $\mathcal E\cap(\supp\nu)^{d+1}$ is compact with positive
probability, and its minimum is at most its conditional average.

Let $\alpha_i(x)$ be barycentric coordinates in the selected simplex.
The first witness is
\[
 V\min\{1,\pos{-\alpha_i(x)}\}.
\]
The pair witness is at least
$V\min\{\pos{\alpha_i(x)},\pos{\alpha_j(x)}\}$:
when both coefficients are positive the tested determinants have opposite
signs and magnitudes $V\alpha_j(x)$ and $V\alpha_i(x)$.
Thus
\[
 \Phi:=\EE\left[
    \sum_i\min\{1,\pos{-\alpha_i(X)}\}
    +\sum_{i<j}\min\{\pos{\alpha_i(X)},\pos{\alpha_j(X)}\}\right]
 \le\frac{N_dD}{q_0v_0}.
\]
Replacement determinants are at most $2^d$, so
$|\alpha_i(x)|\le L_0:=2^d/v_0$ for $x\in\BB$.
Assign $x$ to its largest positive coefficient, with deterministic
tie-breaking, and call this index $r$.
Writing $N_x=\sum_i\pos{-\alpha_i(x)}$ and
$U_x=\sum_{i\ne r}\pos{\alpha_i(x)}$, we have
\[
 N_x\le L_0\sum_i\min\{1,\pos{-\alpha_i(x)}\},
 \qquad
 U_x\le\sum_{i<j}\min\{\pos{\alpha_i(x)},\pos{\alpha_j(x)}\}.
\]
For the second inequality use the pairs containing $r$.
The diameter bound gives
$\|x-w_r\|\le2\sum_{i\ne r}|\alpha_i(x)|=2(N_x+U_x)$.
Therefore
\[
 h\le2L_0\Phi
 \le\frac{2L_0N_d}{q_0v_0}D
 =(d+1)(d+2)8^d b^{-4d}D=QD.
\]
Finally, $\pos{u\cdot x}$ is one-Lipschitz and its expectation is at
least $2b$. The assigned law has expectation at least $2b-h\ge b$
if $QD\le b$. In every direction an anchor therefore has
$u\cdot w_i\ge b$, proving $b\BB\subset T$.
The upper inclusion is automatic.
\end{proof}

\section{Centering controls normalized projection volumes}

Assume now $QD\le b$, select the anchors from Lemma
\ref{lem:assignment}, and put $P=T^\circ$.
It is an enclosing simplex:
\begin{equation}\label{eq:P}
 \BB\subset K\subset P\subset M\BB.
\end{equation}
Each $w_i\in\partial K^\circ$, so $h_K(w_i)=h_P(w_i)=1$.
Let $q_i$ be the vertex of $P$ opposite its facet with normal $w_i$,
and let $\lambda_i$ be the barycentric coordinates of zero in $T$.
They are positive and sum to one.
The cone law of $P$ is
\begin{equation}\label{eq:Plaw}
 \nu_P=\sum_{i=0}^d\lambda_i\delta_{w_i}.
\end{equation}
To see this, the polar images of the facet normals of $P$ are precisely
$w_i$. Its cone law is centered, so its probability weights are the
unique barycentric coordinates of zero.

Write $p_i=\mathbb P(I(X)=i)$ and $c=\sum_i p_iw_i$.
Centering of $\nu$ gives $\|c\|\le h$.
The barycentric coordinate functions of $T$ have the exact formula
\begin{equation}\label{eq:barycentric}
 \alpha_i(x)=\lambda_i(1-q_i\cdot x).
\end{equation}
For $j\ne i$, $q_i\cdot w_j=1$; taking the scalar product of
$\sum_j\lambda_jw_j=0$ with $q_i$ gives
$q_i\cdot w_i=1-1/\lambda_i$.
These values prove \eqref{eq:barycentric} on an affine basis.
Consequently
\begin{equation}\label{eq:weights}
 p_i-\lambda_i=-\lambda_i q_i\cdot c,\qquad
 \sum_i|p_i-\lambda_i|\le M\|c\|\le Mh.
\end{equation}
For any unit $u$, the function $x\mapsto|u\cdot x|$ is one-Lipschitz
and bounded by one on $T$.
Comparing first with the assigned law and then with \eqref{eq:Plaw} gives
\[
 \left|\int|u\cdot x|\,d\nu(x)
       -\int|u\cdot x|\,d\nu_P(x)\right|
 \le (1+M)h.
\]
By \eqref{eq:brightness}, if $\pi_H(u)=|H\mid u^\perp|_{d-1}$, then
\begin{equation}\label{eq:normalized-projections}
 \left|\frac{\pi_K(u)}{|K|}-\frac{\pi_P(u)}{|P|}\right|
 \le\frac d2(1+M)h.
\end{equation}
This is the additional information retained by the integrated assignment.

\section{Scale control and absolute projection deficits}

Since $h_P$ is $M$-Lipschitz and equals one at every anchor,
\[
 0\le z:=\int(h_P(x)-1)\,d\nu(x)\le Mh\le1.
\]
The integral is $V(K[d-1],P)/|K|-1$.
The volume first-variation formula and the Brunn--Minkowski inequality give
Minkowski's first inequality, and hence
\begin{equation}\label{eq:volume}
 \frac{|P|}{|K|}\le(1+z)^d,\qquad
 |P|-|K|\le d2^{d-1}Mh\,|K|.
\end{equation}
The second bound uses $(1+z)^d-1\le d2^{d-1}z$ for $0\le z\le1$.
Explicitly, differentiating
$|K+tP|^{1/d}\ge|K|^{1/d}+t|P|^{1/d}$ at $t=0+$
proves the mixed-volume inequality just used.

Since $\nu_P$ is supported on $\BB$, \eqref{eq:brightness} gives
$\pi_P(u)/|P|\le d/2$.
Using \eqref{eq:normalized-projections} and \eqref{eq:volume},
\begin{align*}
 0\le\pi_P(u)-\pi_K(u)
 &=\frac{\pi_P(u)}{|P|}(|P|-|K|)
    +|K|\left(\frac{\pi_P(u)}{|P|}
              -\frac{\pi_K(u)}{|K|}\right)\\
 &\le\frac d2 |K|
       \bigl[d2^{d-1}M+1+M\bigr]h
 \le Jh.
\end{align*}
The last step uses $|K|\le(2R_0)^d$.
We have thus obtained the uniform absolute projection bound
\begin{equation}\label{eq:projection-deficit}
 \sup_{|u|=1}\{\pi_P(u)-\pi_K(u)\}\le Jh.
\end{equation}
The nonnegativity uses the actual inclusion $K\subset P$.

\section{A cap in one projection}

\begin{lemma}[Projection cap]\label{lem:projection-cap}
If $\BB\subset K\subset P\subset M\BB$, $d\ge2$, and
$\pi_P(u)-\pi_K(u)\le\varepsilon$ in every direction, then
\[
 d_H(K,P)\le(d-1)(M+1)\varepsilon^{1/(d-1)}.
\]
No simplex hypothesis is needed for this lemma.
\end{lemma}

\begin{proof}
Put $s=d_H(K,P)$. If $s=0$ there is nothing to prove.
Otherwise choose a farthest $q\in P$ and its metric projection $k\in K$,
and set $n=(q-k)/s$.
The projection optimality condition gives $n\cdot k=h_K(n)$, so
$n\cdot q-h_K(n)=s$. Also $s\le M$, since $0\in K$.
Choose a unit $u\perp n$ and project orthogonally onto $u^\perp$.
Let $\overline P,\overline K,\overline q$ be the projected objects.
The space has dimension $m=d-1$, $n$ lies in it,
$\overline P$ contains its unit ball and is contained in its ball of
radius $M$, and
$n\cdot\overline q-h_{\overline K}(n)=s$.

With $\tau=s/[2(M+1)]<1$, convexity places the ball
$(1-\tau)\overline q+\tau B_2^m$ inside $\overline P$.
Its lowest $n$-coordinate is at least
\[
 n\cdot q-\tau(M+1)=h_{\overline K}(n)+s/2,
\]
so it is disjoint from $\overline K$.
A cube of side $2\tau/m$ fits in that ball. Therefore
\[
 \varepsilon\ge|\overline P|-|\overline K|
 \ge(2\tau/m)^m
 =\left[\frac{s}{(d-1)(M+1)}\right]^{d-1}.
\]
This proves the claim.
\end{proof}

Equations \eqref{eq:projection-deficit} and Lemma
\ref{lem:projection-cap} yield the simplex-specific surface/assignment
conversion
\begin{equation}\label{eq:newconversion}
 s=d_H(K,P)\le L(Jh)^{1/m}
 \le L[JQ(d+1)e(K)]^{1/m}.
\end{equation}
It is stronger in exponent than the full-dimensional volume-cap conversion.

\section{Retaining every initially chosen maximum simplex}

If $0\le e\le e_\sharp$, then
$QD\le Q(d+1)e\le r_\sharp\le b$ and
\eqref{eq:newconversion} gives $s\le1/(8Md)$.
Because $h_P\ge1$ and $h_K\ge h_P-s$,
$(1-s)P\subset K$.
The initially selected $\Delta$ is maximum, so
\[
 \frac{|\Delta|}{|P|}\ge(1-s)^d\ge1-ds.
\]
Put $\delta=ds\le1/8$.
The barycentric coordinates of the vertices of $\Delta$ in $P$ form
a nonnegative column-stochastic $(d+1)\times(d+1)$ matrix $W$ with
$|\det W|=|\Delta|/|P|\ge1-\delta$.
Each column has norm at most one.
Hadamard's inequality therefore gives squared norm at least
$(1-\delta)^2$ in each column, so each column has an entry at least
$1-2\delta$.

These dominant rows are distinct.
Indeed, the permanent is the probability that independent row selections
from the columns are all distinct and dominates $|\det W|$.
A common dominant row in two columns would give
\[
 |\det W|\le\operatorname{per}W
 \le1-(1-2\delta)^2\le4\delta<1-\delta,
\]
a contradiction.
Match the vertices by these rows.
Each vertex of $\Delta$ is within
$2\delta\operatorname{diam}P\le4Md\,s$ of its matched vertex of $P$.
Taking convex combinations and using $\BB\subset\Delta$ gives
\[
 K\subset P\subset\Delta+4Md\,s\,\BB
          \subset(1+4Md\,s)\Delta.
\]
By \eqref{eq:newconversion},
$E(K,\Delta)\le A_d^\sharp e^{1/m}$.
This also covers $e=0$: then $h=0$, the projection deficits vanish,
$s=0$, and the same matrix argument gives zero excess.
At the threshold,
\[
 A_d^\sharp e_\sharp^{1/m}
   =4MdL(Jr_\sharp)^{1/m}\le\frac12.
\]
For $e>e_\sharp$, \eqref{eq:normalization} gives
$K\subset R_0\Delta$ and hence $E(K,\Delta)\le R_0-1$,
which is covered by $G_d^\sharp$ in \eqref{eq:global}.
Undoing the affine normalization restores the centroid of the arbitrary
initial $S$. Thus the estimate holds for every maximum simplex.

\section{Sharpness and the obstruction to using volume alone}

The pinned truncation source gives the following exact calculation;
the needed parts are derived here.
In unnormalized coordinates let
\[
 P_0=\conv(0,e_1,\ldots,e_d),\qquad
 K_t=\{x_i\ge0,\ t\le\textstyle\sum_i x_i\le1\},\quad 0<t<1.
\]
Put $q=t^{d-1}$ and $c=1/(d-1)!$.
Its facet area vectors and area times support numbers are
\[
 (u_i,b_i)=(-c(1-q)e_i,0),\quad
 (u_+,b_+)=(c\mathbf1,c),\quad
 (u_-,b_-)=(-cq\mathbf1,-ct^d).
\]
The polytope determinant formula for entry005 is
$a=L_{\rm det}/(d|K_t|H_{\rm det})$.
Horizontal $d$-minors and lifted $(d+1)$-minors give respectively
\begin{align*}
 H_{\rm det}
 &=c^d(1-q)^{d-1}[d+1+(d-1)q],\\
 L_{\rm det}
 &=c^{d+1}(1-q)^{d-1}
        [1+(d-1)q-(d-1)t^d-qt^d].
\end{align*}
For the first sum the all-coordinate subset contributes
$c^d(1-q)^d$, while each subset with one of the parallel top/bottom
normals contributes $c^d(1-q)^{d-1}$ or
$c^dq(1-q)^{d-1}$.
For the second sum, the all-coordinate subsets with top or bottom
contribute $c^{d+1}(1-q)^d$ and
$c^{d+1}(1-q)^dt^d$; each of the $d$ remaining subsets contributes
$c^{d+1}(1-q)^{d-1}(q-t^d)$.
Since $|K_t|=(1-t^d)/d!$, subtraction of $1/(d+1)$ gives
\begin{equation}\label{eq:truncation}
 e(K_t)=
 \frac{t^{d-1}[d(d-1)-(d+1)(d-2)t-2t^d]}
 {(d+1)(1-t^d)[d+1+(d-1)t^{d-1}]},
 \quad
 \lim_{t\downarrow0}\frac{e(K_t)}{t^{d-1}}
 =\frac{d(d-1)}{(d+1)^2}.
\end{equation}
Using nonpositive support numbers here is legitimate: this is a facet
determinant calculation, not a probability law. Translation changes
the lifted determinants by a row operation and leaves them unchanged.

The simplex $S_t=\conv(e_1,\ldots,e_d,te_1)$ is maximum.
Indeed a nondegenerate vertex simplex must represent all $d$ coordinate
rays, with one ray doubled. If $k$ of its other selected vertices are
bottom vertices, its lifted determinant has magnitude $(1-t)t^k$.
The maximum is $1-t$.
For arbitrary inscribed vertices, separate affine linearity and the
triangle inequality bound their determinant by that same vertex maximum.
Thus $|S_t|=(1-t)/d!$ is the global maximum.

Writing $s_x=\sum_i x_i$, barycentric coordinates in $S_t$ are
\[
 \beta_0=\frac{1-s_x}{1-t},\qquad
 \beta_1=x_1-t\beta_0,\qquad
 \beta_i=x_i\quad(2\le i\le d).
\]
The minimum of $\beta_1$ on $K_t$ is $-t$, attained at $te_2$;
the other coordinates have minimum zero.
Centroid dilation by factor $1+\epsilon$ is equivalent to
$\beta_i\ge-\epsilon/(d+1)$.
Hence
\[
 E(K_t,S_t)=(d+1)t.
\]
Together with \eqref{eq:truncation}, this forbids every power
$\alpha>1/(d-1)$ in an every-maximum-simplex estimate.
At the endpoint any such universal constant must satisfy
\[
 G_d\ge(d+1)
       \left[\frac{(d+1)^2}{d(d-1)}\right]^{1/(d-1)}.
\]

There is also an exact obstruction to strengthening only the old scalar
mixed-volume step.
For this nested pair,
\[
 d_H(K_t,P_0)=\frac{t}{\sqrt d},\qquad
 \frac{V(K_t[d-1],P_0)}{|K_t|}-1
       =\frac{t^d}{1-t^d}.
\]
The distance maximum occurs at the missing vertex zero; its nearest
point is $(t/d)\mathbf1$.
In the first-variation integral only the bottom facet has a support gap,
namely $t/\sqrt d$, and its area is
$\sqrt d\,t^{d-1}/(d-1)!$. This proves the mixed-volume formula.
Thus an estimate using this scalar alone cannot have a power greater
than $1/d$.
For all sufficiently small $t$, a fixed translation and dilation put
these pairs in a position with a common inner ball and a fixed outer ball;
the obstruction is therefore not caused by the coordinate origin.
In contrast, for $u=(e_1-e_2)/\sqrt2$, Cauchy's formula gives
\[
 \pi_{P_0}(u)-\pi_{K_t}(u)
 =\frac{t^{d-1}}{\sqrt2\,(d-1)!}.
\]
The top and bottom area vectors are perpendicular to $u$, and only
the two coordinate facets contribute. This identifies exactly the
information that the projection argument preserves.

\section{Provenance and verification scope}

The already published $1/d$ manuscript is included separately as
\nolinkurl{sources/prior-integrated-witness/paper.tex}, with its theorem
and constants preserved. Its exact public revision and readable-source
hashes are pinned in \texttt{PROVENANCE.json}. The received historical
fallback differs only by disclosed editorial changes; its original hash
is recorded separately. No inaccessible source is required for replay.

The pinned v3 source was read for the cone-law representation.
The pinned quantitative supplement was read for normalization,
first variation, and the every-maximum conversion.
The actual pinned truncation \texttt{proof.tex} and exact checker were read.
The new centering/brightness and projection-cap arguments above are
self-contained quantitative estimates.

\texttt{check\_exact.py} uses rational arithmetic and checks the barycentric
weight correction, centered truncation cone laws, normalized and absolute
projection identities, mixed-volume and projection scalings, both sets
of explicit threshold constants, and finite near-permutation matrices.
Its moment/witness regressions enumerate actual rational samples.
The inherited truncation checker is replayed in a separate directory.
Finite checks supplement the analytic proof and do not establish its
arbitrary-body or nonatomic assertions.
No inverse-Minkowski stability theorem, including the disputed proof of
arXiv:2603.17726v3, is used.

\section{Formal verification and reproduction}
The principal theorem has a complete Lean proof. The precise scope is:
Actual convex bodies and projection-body deficit in every dimension d $\ge$ 3; upper bound for every maximum-volume inscribed simplex, and failure of every larger exponent using specified actual maximum simplices. Constants are explicit, not optimal.
The public proof source snapshot is fixed at commit
\nolinkurl{02af7db99db789fecb0d5f5f8f37b1a4e40762a0} in \url{https://github.com/mxym/math}.
The proof packages, with their own dependency locks, audit records and
reproduction instructions, are:
\begin{itemize}
\item \nolinkurl{formalizations/sharp-simplex-upper-bound}.
\item \nolinkurl{formalizations/simplex-truncation-sharpness}.
\end{itemize}
The main formal entry points are:
\begin{itemize}
\item \nolinkurl{Entry005.sharpMain}.
\item \nolinkurl{Entry005.truncationSharpness}.
\end{itemize}
Lean is fixed to version 4.34.1, compiler commit
\nolinkurl{5045d0056413266e57c625dcd7c365b10e377c52}; Mathlib is fixed to
\nolinkurl{d13f23b723b8a846827a245b89c10fc7d3f11612}.
The audited mathematical roots depend only on the standard axioms
\texttt{propext}, \texttt{Classical.choice}, and \texttt{Quot.sound}.
The archived audits state their actual compilation and replay scopes;
publication of this typeset version does not imply a new execution of every
previous empty-kernel replay. The companion README maps each printed
principal result to its formal theorem and explains how to repeat verification.

\section*{Research and preparation disclosures}
This research, formalization and manuscript preparation used AI assistance.
The author is responsible for the manuscript. The principal claims are
supported by the readable proofs and the public Lean artifacts specified
above. Formal verification is not external human peer review; no such review
is claimed. No external funding was received. No worldwide priority or
journal acceptance is asserted.

\begin{thebibliography}{9}
\bibitem{entry}
mxym, \emph{Random-determinant rigidity, sharp symmetric cone bounds, and
spectral nonattainment}, entry005 v3, commit
\nolinkurl{31e3d8a37e4a3051a7f5a2535be1c75642e15bcc}.
\url{https://github.com/mxym/math/blob/31e3d8a37e4a3051a7f5a2535be1c75642e15bcc/preprints/005-simplex-product-optimum/v3/paper.md}.
\bibitem{released}
mxym, \emph{Effective simplex rigidity for the projection cone invariant},
Sections 2, 4--6, inspected at commit
\nolinkurl{87868bcb65bf0460a90d6bc2efc3324e46481230}.
\url{https://github.com/mxym/math/blob/87868bcb65bf0460a90d6bc2efc3324e46481230/notes/quantitative-projection-simplex-stability/paper.md}.
\bibitem{truncation}
mxym, \emph{An exact exponent obstruction for lower-end projection-cone
simplex stability}, inspected at the same commit.
\url{https://github.com/mxym/math/blob/87868bcb65bf0460a90d6bc2efc3324e46481230/notes/simplex-truncation-stability/proof.tex}.
\bibitem{classical}
R. Schneider, \emph{Convex Bodies: The Brunn--Minkowski Theory},
second expanded edition, Cambridge University Press, 2014.
Only classical Cauchy, surface area, first-variation and
Brunn--Minkowski identities are used.
\end{thebibliography}
\end{document}
